\documentclass[10pt,twocolumn,letterpaper]{article}

\usepackage{cvpr}
\usepackage{times}
\usepackage{epsfig}
\usepackage{graphicx}
\usepackage{amsmath}
\usepackage{amssymb}
\usepackage{booktabs}
\usepackage{multirow}

\usepackage[breaklinks=true,bookmarks=false]{hyperref}

\cvprfinalcopy

\def\cvprPaperID{****}
\def\httilde{\mbox{\tt\raisebox{-.5ex}{\symbol{126}}}}

\setcounter{page}{1}
\begin{document}

\title{On the Role of Graph Neural Networks in Distributed\\
Reinforcement Learning for Packet Routing}

\author{Miquel Ainaud\\
{\tt\small miquel.ainaud@estudiantat.upc.edu}
\and
Josep Vidal\\
{\tt\small josep.vidal@upc.edu}
}

\maketitle

\begin{abstract}
We implement and study ScaIR, a distributed multi-agent reinforcement
learning algorithm for packet routing that uses a per-node sub-GNN to
produce topology-aware feature vectors as input to a Q-network. Our
starting goal was to extend ScaIR with more expressive aggregation
mechanisms such as graph attention, inspired by similar ideas in recent
work on spatiotemporal routing. Instead, we ended up finding something
more fundamental: the GNN component does not contribute meaningfully to
routing quality. Replacing it entirely with a fixed one-hot node identifier
or a static neighbour connectivity mask achieves equal or better
performance across all topologies tested. The reason is structural: since
$\mathbf{V}_n$ is initialised from a one-hot vector and the GNN has no
access to delays, traffic load, or any physical network information, the
learned feature vectors converge to a transformation of node identity
rather than a representation of topology. This explains both why removing
the GNN does not hurt, and why the model fails badly when a new node is
added to the network---it simply has no way to assign a meaningful
identity to an entity it has never seen. The actual gains over OSPF,
which can reach 40--80\% under congested traffic, come from the
multi-agent DQN and UCB exploration mechanism, not from the GNN. We
discuss what these findings imply for the design of future
topology-aware routing systems.
\end{abstract}

\section{Introduction}

\subsection{The routing problem}

Packet routing is one of the core tasks in any communication network.
When a packet arrives at a router, the router has to decide which
neighbour to forward it to, with the goal of getting it to its
destination as fast as possible. Traditionally this is solved with
protocols like OSPF, which precomputes the shortest path to every
destination using Dijkstra's algorithm \cite{dijkstra1959}. This works
well when traffic is spread evenly across the network, but breaks down
when certain source--destination pairs concentrate a lot of demand:
OSPF will route everything along the same shortest path, which
saturates, while alternative routes that could share the load remain
idle. Minimising average end-to-end packet delay---or equivalently
maximising throughput---under these conditions is a hard combinatorial
problem that depends on the instantaneous state of every queue in the
network, which OSPF simply does not consider.

Over the last few years, reinforcement learning has emerged as a natural
way to tackle this. A router modelled as an RL agent can learn, through
trial and error, to route packets around congested areas by observing
local queue lengths and adjusting its forwarding decisions online. The
combination of RL with graph neural networks (GNNs) has attracted
particular interest, since the network is inherently a graph and GNNs
offer a principled way to aggregate information from neighbours and
encode topology structure \cite{survey_gnn_drl}. Several works have
explored this idea from different angles: some adopt a centralised
software-defined networking approach where a single controller collects
global state and computes routing decisions for the whole network
\cite{upc_routing}, while others look at hybrid or automatic
architecture search to find the right GNN structure for a given network
scenario \cite{autognn}. The consensus is that RL-based routing can
outperform OSPF under congestion, and that incorporating some form of
graph-based feature extraction helps.

The approach we focus on in this work, ScaIR \cite{scair2025}, is a
fully distributed system: each router is an independent agent with its
own sub-GNN and Q-network, and the only communication between agents is
the exchange of feature vectors with direct neighbours. This avoids any
centralised coordinator and scales naturally with network size, which
makes it an attractive design for large real-world deployments.

\subsection{From STRL to ScaIR: the starting point of this work}

One of the papers that motivated this project is STRL
\cite{strl}, which proposes a spatiotemporal RL approach to routing. A
key observation in that work is that the routing problem is
non-Markovian: the delivery time of a packet depends not just on the
current state of the queues but also on recent history, since congestion
builds up and dissipates over time. STRL addresses this by incorporating
temporal features into the state representation, and it also advocates
for a distributed deployment model to improve scalability.

ScaIR \cite{scair2025} shares the distributed philosophy but takes a
different angle on global awareness: instead of temporal features, it
uses a GNN to give each agent a compressed summary of the neighbourhood
topology. The architecture is deliberately simple---a two-layer MLP for
each of the GNN update and output functions, with no attention
mechanisms or learnable aggregation weights. Our initial plan was to
improve this by adding graph attention (GAT-style) to the GNN
aggregation, following ideas from STRL and AutoGNN \cite{autognn}, and
to explore whether temporal features from the packet history could also
help. As we will describe, the experiments pushed us in a different
direction: they revealed that the GNN itself is not doing what the
architecture claims, and that fixing this is a more fundamental
challenge than adding attention on top of it.

\section{The ScaIR Architecture}
\label{sec:arch}

ScaIR~\cite{scair2025} is a fully distributed multi-agent routing
system. Each router in the network is an independent
\emph{intelligent router} (IRr): it holds its own neural network
weights and takes all forwarding decisions autonomously, with no
centralised controller. The only inter-agent communication is the
periodic exchange of compressed feature vectors between one-hop
neighbours. The global optimisation objective is to minimise the mean
packet delivery time across all packets in an episode:
\begin{equation}
  \min \;\frac{1}{P_n}\sum_{p \in P} t_p,
  \label{eq:obj}
\end{equation}
where $P$ is the set of packets and $t_p$ is the total time from
generation to delivery for packet $p$.

\subsection{Per-agent architecture: sub-GNN and Q-network}

Each agent $n$ has two jointly-trained components: a sub-GNN
($N_{Gsub}$) and a Q-network ($N_Q$).

\textbf{Sub-GNN ($N_{Gsub}$).} Each node $n$ maintains a hidden
feature vector $V_n$, initialised as a one-hot encoding of its integer
router identifier $R_n$, zero-padded to feature length $F_l$. The
feature vector is updated iteratively by exchanging representations
with one-hop neighbours:
\begin{align}
  V_n^{(t)} &= f_{w_n}\!\bigl(
    V_n^{(t-1)},\;\{V_j^{(t-1)}\}_{j \in N_n}
  \bigr), \label{eq:gnn} \\
  o_n &= g_{w_n}(V_n^{(K)}), \label{eq:out}
\end{align}
where $f_{w_n}(\cdot)$ updates the hidden representation of node $n$
using the representations of its one-hop neighbours $N_n$,
$g_{w_n}(\cdot)$ maps the final hidden state to the output feature
space, and $K$ is the number of message-passing iterations. Both
$f_{w_n}$ and $g_{w_n}$ are feedforward networks. The paper does not
specify how the neighbour representations $\{V_j\}_{j \in N_n}$ are
aggregated inside $f_{w_n}$; we describe our implementation choices in
Section~\ref{sec:experiments}. After $K$ iterations, the output $o_n$
is a compressed representation of the global network topology as seen
from node $n$.

\textbf{Q-network ($N_Q$).} When processing packet $p$, agent $n$
observes the state
\begin{equation}
  S_n = \{d_p,\; O_n,\; K_n\},
  \label{eq:state}
\end{equation}
where $d_p$ is the destination of the packet (encoded as a one-hot
vector), $O_n$ is the local observation comprising queue lengths and
action history for each outgoing interface, and $K_n = o_n$ is the
sub-GNN output representing global topology characteristics. The
Q-network maps $S_n$ to one Q-value per neighbour; the agent forwards
the packet to the neighbour with the minimum estimated delivery cost.
As the exploration policy, ScaIR uses a $\sigma$-greedy strategy: with
probability $\sigma$ the agent selects a random neighbour from $A_n$;
otherwise it acts greedily and forwards to the neighbour with the
minimum Q-value.

For training, each agent $n$ maintains a replay buffer $D_n$ that
accumulates the transitions $(s_t, a_t, c_t, s_{t+1})$ observed during
routing. Every $L_c$ forwarding decisions a random mini-batch is
sampled from $D_n$ to update the network parameters. Because the
Q-value $Q_n(s, a;\theta_n, g_n)$ is interpreted as the estimated total
delivery time if agent $n$ takes action $a$ in state $s$, the
regression target for each transition is built as the sum of the
immediate cost and the best estimated future cost from the next agent:
\begin{equation}
  y = c_t + \gamma\,\hat{Q}(s_{t+1}, a^*;\hat{\theta}_n, \hat{g}_n),
  \label{eq:target}
\end{equation}
where $c_t = q + h$ is the immediate forwarding cost (queuing time $q$
plus per-hop transmission delay $h$), $\gamma$ is the discount factor,
and $\hat{Q}$ is a periodically frozen copy of the Q-network used as
a target network to stabilise training. The parameters of both
$N_{Gsub}$ (denoted $g_n$) and $N_Q$ (denoted $\theta_n$) are then
updated jointly by minimising the loss
\begin{equation}
  L_t = \bigl(y - Q_n(s_t, a_t;\theta_n, g_n)\bigr)^2.
  \label{eq:loss}
\end{equation}
Differentiating $L_t$ with respect to both $\theta_n$ and $g_n$ and
applying gradient descent simultaneously refines the routing policy and
the GNN feature extractor.

\subsection{Fully learned forwarding}

A defining property of ScaIR is that the routing decision is made
entirely by the Q-network, with no shortest-path solver involved at any
stage. This sets it apart from hybrid frameworks such as
\cite{upc_routing} and \cite{autognn}, where neural networks produce
per-link cost estimates or embeddings that are then fed into a
Dijkstra-based solver to compute the actual forwarding path. In those
approaches the learned component shapes the cost surface, but path
selection remains a graph search. In ScaIR, $N_Q$ directly outputs a
next-hop action for each packet based on the current state $S_n$;
there is no graph search step at inference time.

\subsection{Simulation assumptions}

The environment used in both the original paper and our replication
makes three key simplifying assumptions. First, all links are
identical: there are no per-link weights, bandwidth capacities, or
heterogeneous propagation delays. Second, the only source of variable
latency is queue occupancy. The immediate forwarding cost is $c_t = q
+ h$, where $q$ is the number of packets currently queued ahead on the
chosen interface (each contributing $1$\,ms of delay) and $h = 1$\,ms
is a fixed per-hop transmission time. Third---and most relevant to our
findings---the sub-GNN receives no information about queue lengths,
traffic load, or link properties; the only input flowing into the GNN
is node identity through the one-hot initialisation. Queue and traffic
information enter the Q-network separately through $O_n$. These
simplifications keep the simulation tractable but, as we argue below,
they strongly limit what the GNN can usefully learn to represent.

\subsection{Datasets}

We experiment on four real backbone networks. \textbf{Abilene} (11
nodes, 14 links) and \textbf{GEANT} (23 nodes, 37 links) are the same
topologies used in the original ScaIR paper, with traffic matrices from
the Internet2 and GEANT projects. We additionally include \textbf{BRAIN}
(9 nodes after collapsing the 161-node Berlin backbone to its core
routers, with 8619 one-minute TMs from SNDlib) and \textbf{Germany50}
(50 nodes, 288 five-minute TMs from SNDlib \cite{sndlib}). Testing on
different topology sizes and TM characteristics matters because what
looks like a good result on a small network may not hold at scale, and
the BRAIN and Germany50 TMs already have concentrated demand without any
artificial hot-spot injection.

\section{Experiments, Results, and Analysis}
\label{sec:experiments}

The central question we investigate is whether the sub-GNN in ScaIR
provides meaningful topology information to the Q-network, or whether
simpler static representations of the network achieve equal or better
routing quality. We will show that the GNN output $o_n$ encodes little
more than a transformed node identity: it is not learning a topology
representation, and replacing it with a one-hot vector or a static
connectivity mask delivers the same performance across all topologies
and congestion levels tested.

To build this case, we report the following experiments in order:
\begin{itemize}
  \item \textbf{Shared weights.} Replacing per-agent GNN parameter sets
    with a single shared set across all agents. If the GNN encodes
    topology structure, per-node specialisation should matter; if it
    encodes identity, sharing should be largely free.
  \item \textbf{Attention aggregation.} Replacing mean aggregation inside
    the GNN update with dot-product and learnable attention,
    motivated by related work on graph-based
    routing~\cite{strl,autognn}. Selective neighbour weighting should
    improve the topology encoding if the features carry routing-relevant
    content.
  \item \textbf{Removing the GNN.} Replacing the learned sub-GNN
    entirely with a fixed, non-trained vector (one-hot node identifier
    or neighbour connectivity mask) fed directly into the Q-network.
  \item \textbf{Alternative initialisation.} Seeding $V_n$ with
    structural topology properties (degree, betweenness centrality)
    rather than one-hot identifiers~\cite{init_ref1,init_ref2}, to test
    whether richer initial features allow the GNN to converge to a
    routing-useful encoding.
  \item \textbf{Topology resilience.} Evaluating trained agents after
    structural network changes (node/link additions and removals), to
    examine the operational consequences of identity-based
    representations.
\end{itemize}

\subsection{Datasets and Experimental Setup}
\label{sec:setup}

\subsubsection{Topologies}

We evaluate on four real backbone networks. \textbf{Abilene} (11~nodes,
14~links) is the Internet2 US backbone and \textbf{GEANT} (23~nodes,
37~links) is the pan-European research network; both topologies and
their traffic matrices are used in the original ScaIR
paper~\cite{scair2025}, allowing a direct cross-check of our
implementation. Topology files and TMs were obtained from the GitHub
repository provided in the ScaIR paper.

\textbf{BRAIN} and \textbf{Germany50} come from
SNDlib~\cite{sndlib}. The BRAIN instance has 161 nodes, but 152 are
subnet nodes each connected by a single uplink with no routing choice,
so we collapse to the 9 backbone core routers and aggregate the 8619
one-minute TMs accordingly. Germany50 (50~nodes) is the DFN German
national backbone, with 288 five-minute TMs. These two topologies serve
a dual purpose: their structurally different connectivity patterns and
larger scales let us verify that findings from Abilene and GEANT
generalise, and their real-world TMs carry more concentrated demand,
making them naturally harder for OSPF without artificial injection.

\subsubsection{Traffic generation and distribution ratio}

We follow the ScaIR traffic model. At the start of each episode, $P_n$
packets are initialised at source nodes with source--destination pairs
sampled proportional to a traffic matrix $\mathbf{T}$. A fraction $D_r$
of packets is additionally forced onto a fixed hot-spot pair (the
\emph{distribution ratio}), and we sweep $D_r \in \{0.0, 0.2, 0.4,
0.6, 0.8\}$.

Sweeping $D_r$ traces the transition from near-uniform load, where
OSPF's static shortest-path routing is already close to
optimal~\cite{scair2025}, to severely congested conditions where
hot-spot links saturate and delay compounds across hops. Reinforcement
learning-based routing provides its advantage in the congested regime:
the Q-network learns to distribute load over longer but less-loaded
alternatives that OSPF ignores. The primary baseline throughout is OSPF
with unit link weights.

\subsubsection{Common hyperparameters}

Table~\ref{tab:hyper} lists the parameters common to all experiments.
Values follow the ScaIR paper where specified; deviations are noted per
experiment.

\begin{table}[t]
\centering
\caption{Common hyperparameters for all experiments.}
\label{tab:hyper}
\small
\setlength{\tabcolsep}{4pt}
\begin{tabular}{lll}
\toprule
Parameter & Symbol & Value \\
\midrule
Feature vector length      & $F_l$             & 128 \\
GNN warm-up iterations     & $K$               & 8 \\
GNN periodic update iters  & $I_n$             & 3 \\
GNN update period (ticks)  & $T_{\rm GNN}$     & 10 \\
Q-net learning rate        & $L_r$             & $10^{-3}$ \\
GNN learning rate          & $L_r^{\rm GNN}$   & $10^{-4}$ \\
Replay buffer size         & $M_s$             & 200 \\
Batch size                 & $B_s$             & 64 \\
Learning cycle             & $L_c$             & 10 \\
Packets per episode        & $P_n$             & 50 \\
Packet inter-arrival time  & $G_i$             & 0.5\,ms \\
Discount factor            & $\gamma$          & 1.0 \\
Distribution ratios        & $D_r$             & \{0.0, 0.2, 0.4, 0.6, 0.8\} \\
\bottomrule
\end{tabular}
\end{table}

\subsubsection{Baseline performance}

Before examining the GNN's contribution we verify that our
implementation is correct. Table~\ref{tab:main} shows average delivery
time for the standard ScaIR variant (per-node GNN, mean aggregation,
$\sigma$-greedy exploration, 300~episodes) against OSPF across all four
topologies. Gains under heavy congestion are substantial and consistent
with the original paper: up to 81\% on BRAIN and around 76\% on GEANT
and Germany50 at $D_r = 0.8$. On Abilene with near-uniform Internet2
traffic, OSPF wins at $D_r = 0.0$ because there is no congestion for a
learned policy to exploit. On BRAIN and Germany50, ScaIR wins even at
$D_r = 0.0$ because those TMs already carry concentrated demand.

\begin{table}[h]
\centering
\caption{ScaIR (per-node GNN, $\sigma$-greedy, 300~episodes) vs.\
OSPF at two representative distribution ratios.}
\label{tab:main}
\small
\setlength{\tabcolsep}{4pt}
\begin{tabular}{llccc}
\toprule
Topology & $D_r$ & OSPF (ms) & ScaIR (ms) & Gain \\
\midrule
\multirow{2}{*}{Abilene (11)} & 0.0 & 3.35 & 3.52 & $-5\%$ \\
 & 0.8 & 16.20 & 6.89 & $+57\%$ \\
\midrule
\multirow{2}{*}{BRAIN (9)} & 0.0 & 4.15 & 2.30 & $+44\%$ \\
 & 0.8 & 15.52 & 3.01 & $+81\%$ \\
\midrule
\multirow{2}{*}{GEANT (23)} & 0.0 & 2.92 & 2.74 & $+6\%$ \\
 & 0.8 & 22.87 & 5.42 & $+76\%$ \\
\midrule
\multirow{2}{*}{Germany50 (50)} & 0.0 & 4.89 & 4.13 & $+15\%$ \\
 & 0.8 & 28.03 & 6.77 & $+76\%$ \\
\bottomrule
\end{tabular}
\end{table}

\subsection{Shared Weights}
\label{sec:shared}

In the standard ScaIR architecture each agent $n$ maintains independent
GNN parameters $w_n$. A natural simplification is to share a single
parameter set $w$ across all agents while keeping each node's feature
vector $V_n$ independent. Shared weights reduce the parameter count,
increase effective gradient throughput (each update aggregates signal
from all agents), and simplify deployment. If per-node specialisation of
$w_n$ is important --- for instance because the GNN is learning
node-specific topology structure --- sharing should hurt performance.

We train both variants for 300~episodes with $\sigma$-greedy
exploration ($\sigma = 0.9$, decaying by $0.05$ every 10~episodes to a
floor of $0.1$) on all four topologies. Figure~\ref{fig:variants} shows
the $D_r$ sweep on Germany50, the topology where differences between
variants are most visible; Abilene and BRAIN are qualitatively identical
(max spread $< 0.3$\,ms at any $D_r$).

\begin{figure}[t]
\centering
\includegraphics[width=\linewidth]{results/04_germany50_ucb/comparison_dr_sweep.png}
\caption{$D_r$ sweep on Germany50 (50~nodes, 300~episodes,
$\sigma$-greedy) for per-node and shared-weight variants, with mean and
dot-product attention aggregation. All variants converge to within
0.7\,ms of each other; the shared-weight variants trail marginally at
$D_r = 0.8$.}
\label{fig:variants}
\end{figure}

Across Abilene and BRAIN the best and worst variants differ by at most
0.3\,ms at any $D_r$, well within stochastic variance. On Germany50,
shared-weight variants trail per-node by up to 0.7\,ms at $D_r = 0.8$
--- a consistent but modest gap. Weight sharing is therefore a viable
simplification that costs little in routing quality, which already
suggests that per-node GNN specialisation is not a critical source of
routing information.

\subsection{Attention Aggregation}
\label{sec:attn}

If mean aggregation treats all neighbours as equally informative, one
might expect attention-based aggregation to extract a richer topology
representation by weighting more relevant neighbours more heavily. This
is the intuition behind graph attention networks~\cite{velickovic2018}
and motivates related routing works~\cite{strl,autognn}. We implement
the GNN update using only neighbour feature vectors (following the
paper's formulation):
\begin{equation}
  V_n^{(t)} = \mathrm{agg}\!\bigl(\{f_{w_n}(V_j^{(t-1)})\}_{j \in N_n}\bigr),
  \label{eq:gnn-nbr}
\end{equation}
and substitute two attention variants for the mean aggregation.

\textbf{Dot-product attention.} The weight of neighbour $j$ is:
\begin{equation}
  \alpha_j = \frac{
    \exp\!\bigl(V_n^{(t-1)} \cdot f_{w_n}(V_j^{(t-1)}) / \sqrt{F_l}\bigr)
  }{
    \sum_{k \in N_n}\!\exp\!\bigl(V_n^{(t-1)} \cdot f_{w_n}(V_k^{(t-1)}) / \sqrt{F_l}\bigr)
  },
  \label{eq:dot-attn}
\end{equation}
and the update becomes $V_n^{(t)} = \sum_{j} \alpha_j f_{w_n}(V_j^{(t-1)})$.

\textbf{Learnable attention.} A trainable query matrix
$\mathbf{W}_q \in \mathbb{R}^{F_l \times F_l}$ forms the query:
\begin{equation}
  \alpha_j = \frac{
    \exp\!\bigl((\mathbf{W}_q V_n^{(t-1)}) \cdot f_{w_n}(V_j^{(t-1)}) / \sqrt{F_l}\bigr)
  }{
    \sum_{k \in N_n}\!\exp\!\bigl((\mathbf{W}_q V_n^{(t-1)}) \cdot f_{w_n}(V_k^{(t-1)}) / \sqrt{F_l}\bigr)
  }.
  \label{eq:learn-attn}
\end{equation}

The attention variants appear in Figure~\ref{fig:variants} alongside
the mean variants. On Abilene and BRAIN all four combinations
(mean/dot-attn $\times$ per-node/shared) stay within 0.3\,ms of each
other. On Germany50, per-node dot-product attention reaches 6.77\,ms at
$D_r = 0.8$ versus 7.47\,ms for mean per-node --- the largest gap
observed anywhere, yet still modest.

The reason attention does not help is structural. The score in
Eq.~\eqref{eq:dot-attn} measures the inner product between
$V_n^{(t-1)}$ and the transformed neighbour vector. Since all feature
vectors are initialised from one-hot node identifiers and updated
exclusively from other such vectors, the score measures how similar two
nodes' \emph{identities} are to each other --- not how routing-relevant
a neighbour is. No queue lengths, link loads, or path distances ever
enter the GNN; there is no routing-relevant content for the attention
mechanism to discover. This observation motivates removing the GNN
entirely.

\subsection{Removing the GNN}
\label{sec:nognn}

If the aggregation scheme does not matter, is the GNN contributing
anything at all? We test this directly by replacing the sub-GNN with a
fixed, non-trainable vector of the same length $F_l = 128$, fed as
$K_n$ directly into the Q-network. Only the Q-network trains. We use
two encodings:
\begin{itemize}
  \item \textbf{OneHot}: $V_n = \mathbf{e}_n$, a unit vector at
    position $n$. Pure node identity, no topology information.
  \item \textbf{NbrMask}: $V_n$ is a binary vector with ones at the
    indices of $n$'s direct neighbours. Encodes immediate connectivity
    without any learned transformation.
\end{itemize}

Both encodings are combined with per-node and shared Q-networks, giving
four no-GNN variants. For these experiments we use UCB exploration:
\begin{equation}
  a^* = \arg\min_a \!\left[
    Q_n(S_n, a;\theta_n, g_n) - c\sqrt{\frac{\ln(N+1)}{n_a}}
  \right],
  \label{eq:ucb-full}
\end{equation}
where $N = \sum_a n_a$ is the total forwarding count for the current
destination, $n_a$ is the per-action visit count, and $c = 2.0$.
Unvisited actions are always explored first; counts are
destination-specific and never reset across episodes. The GNN reference
for this section also uses UCB, making the comparison clean.

Figure~\ref{fig:nognn} shows the $D_r$ sweep on Germany50 and
Table~\ref{tab:ablation} gives the numerical results.

\begin{figure}[t]
\centering
\includegraphics[width=\linewidth]{results/05_no_gnn_ger50/ablation_dr_sweep.png}
\caption{No-GNN ablation on Germany50 (UCB, 300~episodes). All
fixed-encoding variants match or outperform ScaIR-GNN at every
$D_r$ value. The GNN adds optimisation complexity without improving
routing quality.}
\label{fig:nognn}
\end{figure}

\begin{table}[t]
\centering
\caption{No-GNN ablation on Germany50 (UCB, 300~episodes). Negative
``vs.\ GNN'' values indicate the no-GNN variant is better.}
\label{tab:ablation}
\small
\setlength{\tabcolsep}{4pt}
\begin{tabular}{lcccc}
\toprule
 & \multicolumn{2}{c}{$D_r = 0.0$} &
   \multicolumn{2}{c}{$D_r = 0.8$} \\
\cmidrule(lr){2-3}\cmidrule(lr){4-5}
Variant & ms & vs.\ GNN & ms & vs.\ GNN \\
\midrule
ScaIR-GNN (per-node) & 4.13 & --- & 6.77 & --- \\
\midrule
OneHot per-node  & 3.93 & $-4.9\%$ & 6.58 & $-2.8\%$ \\
OneHot shared    & 3.94 & $-4.6\%$ & \textbf{6.18} & $-8.7\%$ \\
NbrMask per-node & 3.96 & $-4.1\%$ & 6.63 & $-2.1\%$ \\
NbrMask shared   & \textbf{3.92} & $-5.1\%$ & 6.37 & $-5.9\%$ \\
\midrule
OSPF             & 4.89 & $+18.5\%$ & 28.03 & $+314\%$ \\
\bottomrule
\end{tabular}
\end{table}

Every no-GNN variant matches or outperforms ScaIR-GNN at both
congestion levels. The best (OneHot shared-Q) beats it by 8.7\% at
$D_r = 0.8$; the same holds on BRAIN (9~nodes), where no-GNN variants
differ from the GNN reference by at most 0.18\,ms. The GNN not only
fails to help --- it slightly hurts performance within a 300-episode
budget, most likely because the coupled optimisation of $f_{w_n}$,
$g_{w_n}$, and $N_Q$ introduces gradient instability that a fixed
encoding avoids. The actual routing signal is in the local observations
$O_n$ (queue lengths and action history), which bypass the GNN.

\subsection{Alternative Initialisation}
\label{sec:altinit}

A natural objection is that the issue lies not in the GNN architecture
but in the one-hot initialisation: a $V_n$ seeded from a pure identity
vector will inevitably converge to an identity-like encoding. Perhaps
seeding $V_n$ with structural topology properties would allow the GNN to
learn routing-relevant representations, as explored in graph
representation learning works~\cite{init_ref1,init_ref2}.

We test this by initialising $V_n$ with two topology-derived scalars
(zero-padded to $F_l$):
\begin{itemize}
  \item \textbf{Degree}: $V_n^{(0)}[0] = \deg(n) / \max_k \deg(k)$.
  \item \textbf{Betweenness}: additionally $V_n^{(0)}[1] = \mathrm{BC}(n)$
    (normalised betweenness centrality).
\end{itemize}

Each initialisation is combined with the three aggregation schemes from
Section~\ref{sec:attn} (mean, dot-attention, learnable~attention),
giving six GNN variants. We also include two no-GNN baselines that feed
the topology vector directly to the Q-network without message passing
(\texttt{deg\_fixed}, \texttt{bet\_fixed}), for nine variants in total.
All are trained for 200~episodes with UCB on all four topologies at
$D_r \in \{0.0, 0.4, 0.8\}$.

Figures~\ref{fig:init-ger50} and~\ref{fig:init-abi} show the results on
Germany50 and Abilene; Table~\ref{tab:initspread} summarises the spread
between best and worst variant per topology.

\begin{figure}[t]
\centering
\includegraphics[width=\linewidth]{results/09_topo_init/germany50_dr_sweep.png}
\caption{Alternative initialisation sweep on Germany50 (UCB,
200~episodes). All nine variants, including topology-derived
initialisations and attention, land within 1\,ms of each other.}
\label{fig:init-ger50}
\end{figure}

\begin{figure}[t]
\centering
\includegraphics[width=\linewidth]{results/09_topo_init/abilene_dr_sweep.png}
\caption{Alternative initialisation sweep on Abilene (UCB,
200~episodes).}
\label{fig:init-abi}
\end{figure}

\begin{table}[t]
\centering
\caption{Spread (best minus worst delivery time in ms) across all nine
variants per topology and $D_r$.}
\label{tab:initspread}
\small
\setlength{\tabcolsep}{4pt}
\begin{tabular}{lccc}
\toprule
Topology  & $D_r = 0.0$ & $D_r = 0.4$ & $D_r = 0.8$ \\
\midrule
Abilene   & 0.055 & 0.118 & 0.414 \\
GEANT     & 0.086 & 0.121 & 0.568 \\
BRAIN     & 0.088 & 0.097 & 0.150 \\
Germany50 & 0.352 & 0.473 & 1.037 \\
\bottomrule
\end{tabular}
\end{table}

The spread between best and worst variant is 0.05--1.04\,ms across all
settings. On Germany50 at $D_r = 0.8$ --- where a better topology
encoding might plausibly matter --- the best performer is
\texttt{deg\_fixed} (6.63\,ms), a no-GNN variant that feeds the degree
vector directly to the Q-network without any message passing. The worst
is the learnable-attention degree variant (7.66\,ms); the most complex
model is the worst performer, likely because 200~episodes is
insufficient for $\mathbf{W}_q$ to converge. Every topology-derived
initialisation produces results indistinguishable from the one-hot
baseline within run-to-run noise. The conclusion holds: within a
short training budget, the GNN's topology encoding is a free parameter
and what the Q-network routes on is the local congestion signal in $O_n$.

\subsection{Topology Resilience}
\label{sec:resilience}

So far the finding that $V_n$ converges to an identity-like encoding is
a neutral result: it does not help, but it does not actively hurt
either. The picture changes when the network topology itself is
modified. Any change that introduces a previously unseen node identifier
--- specifically, adding a new router --- will produce an
out-of-distribution input that a trained network cannot handle
gracefully.

We evaluate this using the nine variants from
Section~\ref{sec:altinit}, trained for 200~episodes on Abilene at
$D_r = 0.6$. We apply two mutations: inserting a new router (node~11,
connected to existing nodes~0 and~8), and removing link~3--6.
Performance is measured zero-shot immediately after the change and
after 50 online adaptation episodes.

\textbf{Link removal is handled gracefully.} After removing link~3--6
all nine variants achieve 4.18--4.92\,ms zero-shot, already 36--43\%
better than OSPF on the pre-change topology. The trained routing
policies do not rely critically on any single link.

\textbf{Node insertion exposes the identity encoding.} Adding node~11
causes the three mean-aggregation GNN variants to fail
catastrophically: 162--297\,ms zero-shot, compared to a pre-change
baseline of $\approx 4$\,ms. During the $K = 8$ GNN warm-up iterations
at the start of each episode, node~11's uninitialised feature vector
$V_{11}$ enters the neighbourhood average of nodes~0 and~8, corrupting
their converged $V_n$ representations before the Q-network can act.
Attention variants are considerably more robust: the
degree+learnable-attention variant achieves 5.57\,ms zero-shot, because
it can downweight the contribution of the unfamiliar neighbour vector.
Figure~\ref{fig:resilience} shows zero-shot and post-adaptation results
for all variants.

\begin{figure}[t]
\centering
\includegraphics[width=\linewidth]{results/09_topo_init/adaptation_plot.png}
\caption{Topology adaptation on Abilene ($D_r = 0.6$): zero-shot and
after 50 adaptation episodes for node insertion and link removal. Mean-GNN
variants fail catastrophically on node insertion; attention and
fixed-encoding variants are robust. All variants recover within 50
episodes.}
\label{fig:resilience}
\end{figure}

All variants recover fully within 50 online adaptation episodes
(4.0--4.7\,ms). The key takeaway is not specific to any single variant:
the brittleness under node insertion is a consequence of
parameterising the model by node identity. Any architecture relying on
identity-based feature vectors will fail on unseen node IDs unless it
can attenuate their influence, as the attention mechanism does here.

\section{Conclusions}

The goal of this project was to improve ScaIR's GNN component. What we
found instead was that the component does not do what its design
implies, and that understanding why leads to a more fundamental
conclusion about this class of models.

The sub-GNN in ScaIR does not learn a useful topology representation.
Because it is initialised from one-hot node identifiers and updated
exclusively from other node-identifier vectors, it converges to an
encoding that is functionally equivalent to a static identity label.
Replacing it with a fixed one-hot or a static neighbour-connectivity
mask achieves the same or strictly better routing performance across all
four topologies tested, within a 300-episode training budget. Replacing
the mean aggregation with dot-product or learnable attention yields
essentially the same result, because attention computed over
node-identity vectors carries no routing-relevant signal. Seeding the
GNN with structural topology properties (degree, betweenness) does not
change the picture either: within the training budget, the Q-network
adapts equally well to any fixed vector it receives, and the GNN's
learned updates add no benefit.

The routing gains over OSPF that ScaIR does achieve---and they are
genuine, reaching up to 81\% reduction in delivery time under heavy
congestion---come from the multi-agent DQN's ability to respond to
local queue observations and from UCB exploration. The ScaIR
paper~\cite{scair2025} argues that the GNN plays a key role based on
feature-length ablations, but those experiments do not isolate the GNN:
reducing $F_l$ also shrinks the destination one-hot and action-history
blocks in the Q-network input, so the observed sensitivity is better
explained by the Q-network needing a minimum input dimensionality to
represent its routing policy. A fixed one-hot of sufficient length works
just as well.

The identity encoding does have a direct operational cost at one
specific boundary condition: inserting a new router. When a new node
joins the network, its one-hot feature is unknown to the trained
$f_{w_n}$, and mean-aggregation GNN variants propagate this unknown
vector into the representations of its neighbours, causing catastrophic
performance degradation. Attention variants are more robust because they
can downweight the unfamiliar input. All variants recover within 50
online episodes, but the zero-shot failure of mean-GNN variants is a
genuine deployment concern.

The most likely reason the GNN does not contribute is the simulation
environment itself. All links share the same fixed 1\,ms transmission
time and there are no bandwidth constraints or propagation delays---the
only source of variable latency is queue length. In this setting, a
node's feature vector could at most encode which other nodes are its
neighbours, which is exactly what a one-hot already encodes. In a
richer environment with heterogeneous link capacities and geographic
propagation delays, the GNN would have physically meaningful
information to aggregate, and a learned representation could plausibly
outperform a fixed one. We believe that is the right direction for
future work, alongside replacing the one-hot node encodings throughout
the architecture with permutation-invariant structural features---a
change that would allow the model to generalise across topologies rather
than memorising node identities.

{\small
\bibliographystyle{ieee_fullname}
\begin{thebibliography}{10}

\bibitem{dijkstra1959}
E.~W. Dijkstra.
\newblock A note on two problems in connexion with graphs.
\newblock {\em Numerische Mathematik}, 1:269--271, 1959.

\bibitem{survey_gnn_drl}
X.~Shi, T.~Zhang, and others.
\newblock A survey on graph neural networks and deep reinforcement learning
  for network optimization.
\newblock {\em IEEE Communications Surveys \& Tutorials}, 2023.

\bibitem{upc_routing}
J.~Su{\'a}rez-Varela et~al.
\newblock Exploring a {GNN}-based routing optimization use case.
\newblock In {\em IEEE/IFIP Network Operations and Management Symposium
  (NOMS)}, 2022.

\bibitem{autognn}
R.~Almasan, J.~Su{\'a}rez-Varela, P.~Barlet-Ros, and A.~Cabellos-Aparicio.
\newblock {AutoGNN}: Autograph neural network architecture for network management.
\newblock {\em arXiv preprint arXiv:2201.10286}, 2022.

\bibitem{strl}
H.~Yu, J.~Zheng, B.~Li, and J.~Ren.
\newblock Spatio-temporal graph reinforcement learning for adaptive traffic
  routing.
\newblock In {\em IEEE International Conference on Communications (ICC)}, 2022.

\bibitem{scair2025}
J.~Zhang, J.~Guan, K.~Liu, Y.~Hu, A.~Shen, and Y.~Ma.
\newblock {ScaIR}: Scalable intelligent routing based on distributed graph
  reinforcement learning.
\newblock {\em Computer Networks}, 257:110915, 2025.

\bibitem{you2020}
X.~You, X.~Li, Y.~Xu, H.~Feng, J.~Zhao, and H.~Yan.
\newblock Toward packet routing with fully distributed multiagent deep
  reinforcement learning.
\newblock {\em IEEE Transactions on Systems, Man, and Cybernetics},
  52(2):855--868, 2020.

\bibitem{boyan1994}
J.~A. Boyan and M.~L. Littman.
\newblock Packet routing in dynamically changing networks: A reinforcement
  learning approach.
\newblock In {\em Advances in Neural Information Processing Systems (NIPS)},
  pages 671--678, 1994.

\bibitem{knight2011}
S.~Knight, H.~X. Nguyen, N.~Falkner, R.~Bowden, and M.~Roughan.
\newblock The internet topology zoo.
\newblock {\em IEEE Journal on Selected Areas in Communications},
  29(9):1765--1775, 2011.

\bibitem{sndlib}
S.~Orlowski, R.~Wess{\"a}ly, M.~Pi{\'o}ro, and A.~Tomaszewski.
\newblock {SNDlib} 1.0---survivable network design library.
\newblock {\em Networks}, 55(3):276--286, 2010.

\bibitem{velickovic2018}
P.~Veli{\v{c}}kovi{\'c}, G.~Cucurull, A.~Casanova, A.~Romero, P.~Li{\`o},
  and Y.~Bengio.
\newblock Graph attention networks.
\newblock In {\em International Conference on Learning Representations
  (ICLR)}, 2018.

\bibitem{init_ref1}
[REF1 --- to be filled in by author].

\bibitem{init_ref2}
[REF2 --- to be filled in by author].

\end{thebibliography}
}

\end{document}
